% -----------------------------------------------------------------------------
% -----------------------------------------------------------------------------

\begin{exbox}[谐振子传播核]
自由粒子核只含动能；真正检验“二次作用量可以精确高斯积分”的下一步，是一维谐振子。它的拉格朗日量为
\begin{equation}
 \boxed{L=\frac{m_{\rm qm}}2\dot\xi^2
 -\frac{m_{\rm qm}\omega_0^2}{2}\xi^2,}
 \label{eq:harmonic-lagrangian-dp68}
\end{equation}
其中 $m_{\rm qm}>0$ 是粒子质量，$\omega_0>0$ 是谐振角频率。对 $T=t_f-t_i>0$，固定端点 $\xi(0)=\xi_i$、$\xi(T)=\xi_f$ 的经典轨道为
\begin{equation}
 \boxed{
 \xi_{\rm cl}(t)
 =\frac{\xi_i\sin[\omega_0(T-t)]
 +\xi_f\sin(\omega_0t)}{\sin(\omega_0T)}.}
 \label{eq:harmonic-classical-path-dp64}
\end{equation}
先取 $0<T<\pi/\omega_0$ 避开第一焦散点。沿这条经典轨道的作用量是
\begin{equation}
 \boxed{
 S_{\rm cl}
 =\frac{m_{\rm qm}\omega_0}{2\sin(\omega_0T)}
 \left[
 (\xi_f^2+\xi_i^2)\cos(\omega_0T)-2\xi_f\xi_i
 \right].}
 \label{eq:harmonic-classical-action-dp64}
\end{equation}
由于作用量对路径是二次泛函，任意路径都可写成经典轨道加零端点涨落，而且两部分在作用量中精确解耦。最终传播核为
\begin{equation}
 \boxed{
 K_{\rm ho}(\xi_f,T;\xi_i,0)
 =\sqrt{\frac{m_{\rm qm}\omega_0}
 {2\pi\ii\hbar\sin(\omega_0T)}}
 \exp\!\left[\frac{\ii}{\hbar}S_{\rm cl}(\xi_f,\xi_i;T)\right].}
 \label{eq:harmonic-kernel-dp64}
\end{equation}
这个结果比自由粒子多出的不是某个经验相位，而是两个可以分别追溯的结构：指数来自满足端点条件的\emph{经典}谐振轨道作用量，前因子来自端点固定的量子涨落高斯积分。

先检查 $\omega_0T\ll1$。展开
\begin{equation*}
 \sin(\omega_0T)=\omega_0T+O((\omega_0T)^3),
 \qquad
 \cos(\omega_0T)=1-\frac{\omega_0^2T^2}{2}+\cdots,
\end{equation*}
得到
\begin{equation*}
 S_{\rm cl}
 =\frac{m_{\rm qm}(\xi_f-\xi_i)^2}{2T}
 +O(T),
\end{equation*}
而前因子趋于
\begin{equation*}
 \sqrt{\frac{m_{\rm qm}}{2\pi\ii\hbar T}}.
\end{equation*}
所以\eqnrefs{eq:harmonic-kernel-dp64} 连续回到自由粒子核~\eqnrefs{eq:free-particle-kernel-dp61}。这说明短时间内势能尚未来得及改变主导的局域扩散尺度。

另一个有用的特殊时刻是四分之一周期。取四分之一周期
\begin{equation*}
 T=\frac{\pi}{2\omega_0},
\end{equation*}
则 $\sin(\omega_0T)=1$、$\cos(\omega_0T)=0$，于是
\begin{equation*}
 K_{\rm ho}
 =\sqrt{\frac{m_{\rm qm}\omega_0}{2\pi\ii\hbar}}
 \exp\!\left(-\frac{\ii m_{\rm qm}\omega_0}{\hbar}\xi_f\xi_i\right).
\end{equation*}
这个核除了尺度因子外正是 Fourier 核：谐振子演化四分之一周期把初始坐标分布旋转到共轭动量分布。用相空间语言说，谐振子把 $(\xi,p)$ 作角度 $\omega_0T$ 的刚性旋转；时间透镜、相空间剪切与电子波包聚焦都具有同一种“传播即相空间线性变换”的结构。

再看 $T\to\pi/\omega_0$。此时 $\sin(\omega_0T)\to0$，坐标表示的前因子发散；这不是演化算符失去幺正性，而是所有经典轨道在半周期后把 $\xi_i$ 聚焦到 $-\xi_i$，坐标核趋向带相位的 $\delta(\xi_f+\xi_i)$。跨过这些焦散点以后，平方根分支会积累 Maslov 相位。因此闭\eqnrefs{eq:harmonic-kernel-dp64} 的局部分支由 $T\to0^+$ 的自由核归一化固定，较长时间则通过核的组合律连续延拓。


\medskip

\eqnrefs{eq:harmonic-lagrangian-dp68} 给出二次势中的经典动力学。这个推导只解决一个问题：在给定初末位置后，哪一条经典轨道满足 Euler--Lagrange 方程。端点条件决定两个积分常数，因此不需要任何量子假设。

Euler--Lagrange 方程为
\begin{equation*}
 \ddot\xi_{\rm cl}+\omega_0^2\xi_{\rm cl}=0.
\end{equation*}
一般解写成
\begin{equation*}
 \xi_{\rm cl}(t)=A\cos(\omega_0t)+B\sin(\omega_0t).
\end{equation*}
由 $\xi_{\rm cl}(0)=\xi_i$ 得 $A=\xi_i$；再用 $\xi_{\rm cl}(T)=\xi_f$，
\begin{equation*}
 B=\frac{\xi_f-\xi_i\cos(\omega_0T)}{\sin(\omega_0T)}.
\end{equation*}
代回一般解并利用
\begin{equation*}
 \sin[\omega_0(T-t)]
 =\sin(\omega_0T)\cos(\omega_0t)
 -\cos(\omega_0T)\sin(\omega_0t),
\end{equation*}
即可整理成\eqnrefs{eq:harmonic-classical-path-dp64}。


\medskip

\eqnrefs{eq:harmonic-classical-path-dp64} 已经固定了满足两端边界条件的经典轨道。沿满足运动方程的轨道，谐振子拉格朗日量化成全导数，因此经典作用量由两个端点的 $\xi_{\rm cl}\dot\xi_{\rm cl}$ 直接确定。

利用运动方程，
\begin{align*}
 \frac{\dd}{\dd t}(\xi_{\rm cl}\dot\xi_{\rm cl})
 &=\dot\xi_{\rm cl}^2+\xi_{\rm cl}\ddot\xi_{\rm cl}\\
 &=\dot\xi_{\rm cl}^2-\omega_0^2\xi_{\rm cl}^2.
\end{align*}
所以
\begin{align*}
 S_{\rm cl}
 &=\frac{m_{\rm qm}}2\int_0^T\dd t
 (\dot\xi_{\rm cl}^2-\omega_0^2\xi_{\rm cl}^2)\\
 &=\frac{m_{\rm qm}}2[\xi_{\rm cl}\dot\xi_{\rm cl}]_0^T.
\end{align*}
由\eqnrefs{eq:harmonic-classical-path-dp64}，
\begin{align*}
 \dot\xi_{\rm cl}(0)
 &=\frac{\omega_0}{\sin(\omega_0T)}
 [\xi_f-\xi_i\cos(\omega_0T)],\\
 \dot\xi_{\rm cl}(T)
 &=\frac{\omega_0}{\sin(\omega_0T)}
 [\xi_f\cos(\omega_0T)-\xi_i].
\end{align*}
将 $\xi_{\rm cl}(0)=\xi_i$、$\xi_{\rm cl}(T)=\xi_f$ 与这两个端点速度代入上式，合并同类项后得到\eqnrefs{eq:harmonic-classical-action-dp64}。


\medskip

经典作用量只决定传播核的端点相位；还需要确定所有零端点量子涨落给出的归一化前因子。因为\eqnrefs{eq:harmonic-lagrangian-dp68} 对路径是二次的，将路径写成经典轨道与零端点涨落之和后，交叉项由经典运动方程精确消失，剩余涨落积分只依赖传播时间而不依赖两个端点。

写
\begin{equation*}
 \xi(t)=\xi_{\rm cl}(t)+\delta\xi(t),
 \qquad \delta\xi(0)=\delta\xi(T)=0.
\end{equation*}
作用量展开为
\begin{align*}
 S[\xi]
 =&\,S_{\rm cl}
 +m_{\rm qm}\int_0^T\dd t
 [\dot\xi_{\rm cl}\,\delta\dot\xi-\omega_0^2\xi_{\rm cl}\,\delta\xi]\\
 &+\frac{m_{\rm qm}}2\int_0^T\dd t
 ((\delta\dot\xi)^2-\omega_0^2(\delta\xi)^2).
\end{align*}
中间的线性项分部积分后为
\begin{equation*}
 -m_{\rm qm}\int_0^T\dd t
 (\ddot\xi_{\rm cl}+\omega_0^2\xi_{\rm cl})\,\delta\xi=0,
\end{equation*}
故
\begin{equation*}
 K_{\rm ho}=N(T)\exp\!\left(\frac{\ii}{\hbar}S_{\rm cl}\right).
\end{equation*}
把这一形式代入传播核满足的 Schr\"odinger 方程，端点相关的 $O(\hbar^0)$ 项自动满足 Hamilton--Jacobi 方程，剩余项给
\begin{equation*}
 \frac{\dot N}{N}
 =-\frac{1}{2m_{\rm qm}}
 \frac{\partial^2S_{\rm cl}}{\partial\xi_f^2}
 =-\frac{\omega_0}{2}\cot(\omega_0T).
\end{equation*}
积分得到
\begin{equation*}
 N(T)=\frac{C}{\sqrt{\sin(\omega_0T)}}.
\end{equation*}
短时间极限必须恢复自由粒子核，而 $\sin(\omega_0T)\sim\omega_0T$，因此
\begin{equation*}
 C=\sqrt{\frac{m_{\rm qm}\omega_0}{2\pi\ii\hbar}}.
\end{equation*}
代回 $K_{\rm ho}=N\exp(\ii S_{\rm cl}/\hbar)$ 就得到\eqnrefs{eq:harmonic-kernel-dp64}。
\end{exbox}






